%!TEX encoding = UTF8
%!TEX root = 0-notes.tex

\chapter{Suites numériques}
\label{chap:2}

\section{Comparaison asymptotiques}

\subsection{Petit o}

\subsection{Grand O}

\subsection{Équivalence}

\section{Convergence}

\subsection{Suite de Cauchy}

% application : série converge ssi la queue tend vers 0
% 
% exemple : dans Q qui sort de Q

\subsection{Théorème de Bolzano-Weierstrass}

\subsection{Limites sup et inf et valeur d'adhérence}

\section{Suites récurrentes et points fixes}

% à mettre au chapitre fonctions ?
% ou alors ça motive la fonction k-lipschitzienne
% THÉOREME ACCROISSEMENTS FINIS !!

\section*{Exercices supplémentaires}
\addcontentsline{toc}{section}{Exercices supplémentaires}

\exe{,difficulty=0}{
	Considérons la fonction réelle $f(x) = \frac12x + \frac1{x}$ sur $\D = [1 ; 1,5]$ ainsi que la suite récurrente
		\[ \begin{cases*} u_0 = 1, \\ u_{n+1} = f(u_n). \end{cases*}. \]
	\begin{enumerate}
		\item
		Montrer que l'ensemble $\D$ est stable par $f$.
		\item
		Calculer $f'$ et vérifier que $|f'(x)| \leq c < 1$ pour tout $x\in\D$ 
		pour un certain $c$ à déterminer.
		\item
		Conclure à l'aide du théorème des accroissements finis que, pour tout $x, y \in \D$,
			\[ |f(x) - f(y)| < c |x-y| \]
		et donc que, pour tout $m<n \in\N$,
			\[ |u_n - u_m| < c^m |u_{n-m} - u_0| \leq 3 c^m. \]
		\item
		Conclure que $(u_n)$ est de Cauchy et qu'elle converge vers une limite $l\in\R$.
		\item
		Calculer cette limite $l$.
	\end{enumerate}
}{exe:sqrt2-contractante}{
	todo
}

%\section*{Solutions}
%\addcontentsline{toc}{section}{Solutions}
%
%\shipoutAnswer



